\documentclass[12pt,a4paper,oneside]{article}
\usepackage[brazil]{babel}
\usepackage[utf8]{inputenc}
\usepackage{olymp}
\usepackage{color}
\usepackage[dvipsnames]{xcolor}
\usepackage{listings}

\setcounter{problem}{1}

\begin{document}

\begin{problem}{Números de Catalan}{catalan}{2}
 
Os números de Catalan formam uma sequência de números naturais que
aparecem em vários problemas de contagem, geralmente envolvendo objetos
definidos recursivamente. Os primeiros números da série, para $n = 0, 1, 2, 3, \dots$ são $1, 1, 2, 5, 14, 42, 132, 429, 1430, 4862, 16796, 58786, \dots$

O $n$-ésimo número da série é denotado por $C_n$ e pode ser calculado de diversas formas. Uma delas é a seguinte:

\begin{equation}
C_n = \frac{4n-2}{n+1}C_{n-1}\notag
\end{equation}

com $C_0=1$.

A sequência foi descoberta por Euler quando contava o número de maneiras que um polígono podia ser dividido em triângulos. O nome é uma homenagem ao matemático Eugène Catalan, que descobriu a conexão da sequência com
expressões entre parênteses.

Implemente uma função recursiva para calcular o $n$-ésimo termo da série.

%\Notes
%
%Observações, se tiver.

\InputFile

A entrada contém vários casos de teste. Cada caso de teste é uma linha contendo um número inteiro $n$. Restrição: $0 \leq n \leq 15$.

 O valor $-1$ encerra a lista de casos de teste e não deve ser processado.

\OutputFile

Seu programa deve gerar uma linha de saída para cada caso de teste, contendo o valor de $C_n$.

%\Notes
%Nesta questão, seu programa deve usar a função \texttt{fatorial} dada %acima exatamente como está, não a modifique! Sua
%tarefa é apenas criar o programa com a função \texttt{main}.

\Example

\begin{example}
\exmp{
1
0
3
5
10
-1
}{
1
1
5
42
16796
}%
\end{example}

%Comentários sobre o exemplo, se necessário.

\end{problem}

\end{document}
